\documentclass[10pt,a4paper]{amsart}
\usepackage{amsmath,amssymb,amsthm,graphicx}
% Retained minimal preamble for the Stacks Project chapter more-algebra.tex (More on Algebra).
% The chapter file pinned-original-source.tex carries only \input{preamble} on its first line, so no
% preamble text lives inside the chapter itself.  The Stacks Project uses one single shared preamble for
% every chapter; the snapshot stages it identically as sources/stacks-chow/preamble.tex and
% sources/stacks-etale-cohomology/preamble.tex.  The block below is copied mechanically, line for line,
% from sources/stacks-chow/preamble.tex lines 14-18 (the Xy-pic setup) and 154-217 (the native theorem
% declarations together with the mathematical macros).  Nothing was added, removed or redefined; the
% wrapper class and the amsmath/amssymb/amsthm/graphicx base packages are supplied by the assembler.
% Chapter: more-algebra.tex -- More on Algebra.
% For commutative diagrams we use Xy-pic
\usepackage[all]{xy}

% We use 2cell for 2-commutative diagrams.
\xyoption{2cell}

% Theorem environments.
%
\theoremstyle{plain}
\newtheorem{theorem}[subsection]{Theorem}
\newtheorem{proposition}[subsection]{Proposition}
\newtheorem{lemma}[subsection]{Lemma}

\theoremstyle{definition}
\newtheorem{definition}[subsection]{Definition}
\newtheorem{example}[subsection]{Example}
\newtheorem{exercise}[subsection]{Exercise}
\newtheorem{situation}[subsection]{Situation}

\theoremstyle{remark}
\newtheorem{remark}[subsection]{Remark}
\newtheorem{remarks}[subsection]{Remarks}

\numberwithin{equation}{subsection}

% Macros
%
\def\lim{\mathop{\mathrm{lim}}\nolimits}
\def\colim{\mathop{\mathrm{colim}}\nolimits}
\def\Spec{\mathop{\mathrm{Spec}}}
\def\Hom{\mathop{\mathrm{Hom}}\nolimits}
\def\Ext{\mathop{\mathrm{Ext}}\nolimits}
\def\SheafHom{\mathop{\mathcal{H}\!\mathit{om}}\nolimits}
\def\SheafExt{\mathop{\mathcal{E}\!\mathit{xt}}\nolimits}
\def\Sch{\mathit{Sch}}
\def\Mor{\mathop{\mathrm{Mor}}\nolimits}
\def\Ob{\mathop{\mathrm{Ob}}\nolimits}
\def\Sh{\mathop{\mathit{Sh}}\nolimits}
\def\NL{\mathop{N\!L}\nolimits}
\def\CH{\mathop{\mathrm{CH}}\nolimits}
\def\proetale{{pro\text{-}\acute{e}tale}}
\def\etale{{\acute{e}tale}}
\def\QCoh{\mathit{QCoh}}
\def\Ker{\mathop{\mathrm{Ker}}}
\def\Im{\mathop{\mathrm{Im}}}
\def\Coker{\mathop{\mathrm{Coker}}}
\def\Coim{\mathop{\mathrm{Coim}}}

% Boxtimes
%
\DeclareMathSymbol{\boxtimes}{\mathbin}{AMSa}{"02}

%
% Macros for moduli stacks/spaces
%
\def\QCohstack{\mathcal{QC}\!\mathit{oh}}
\def\Cohstack{\mathcal{C}\!\mathit{oh}}
\def\Spacesstack{\mathcal{S}\!\mathit{paces}}
\def\Quotfunctor{\mathrm{Quot}}
\def\Hilbfunctor{\mathrm{Hilb}}
\def\Curvesstack{\mathcal{C}\!\mathit{urves}}
\def\Polarizedstack{\mathcal{P}\!\mathit{olarized}}
\def\Complexesstack{\mathcal{C}\!\mathit{omplexes}}
% \Pic is the operator that assigns to X its picard group, usage \Pic(X)
% \Picardstack_{X/B} denotes the Picard stack of X over B
% \Picardfunctor_{X/B} denotes the Picard functor of X over B
\def\Pic{\mathop{\mathrm{Pic}}\nolimits}
\def\Picardstack{\mathcal{P}\!\mathit{ic}}
\def\Picardfunctor{\mathrm{Pic}}
\def\Deformationcategory{\mathcal{D}\!\mathit{ef}}
\usepackage{cleveref}
\crefname{originallemma}{original lemma}{original lemmas}
\Crefname{originallemma}{Original lemma}{Original lemmas}
\makeatletter
\newlabel{more-algebra-lemma-derivation-extends}{{07PE}{source}{}{sourcepointer-1}{}}
\newlabel{more-algebra-lemma-derivation-extends@cref}{{[originallemma][0][]07PE}{[1][0][]source}{}{}{}}
\makeatother
\title{Extending a derivation to a completion, to a localisation, and to a finite type extension after multiplying by a power}
\author{The Stacks Project Authors}
\date{stacks-project:more-algebra.tex:local-snapshot-20261012; GNU FDL 1.2 or later}
\begin{document}
\maketitle
\setcounter{subsection}{1}
\noindent\textit{Editorial note.} All delivered mathematical text is the author's own, copied from the pinned Stacks Project chapter file more-algebra.tex; no statement or proof step was rewritten, reordered or removed.\par\medskip
\noindent\textit{Editorial note.} The chapter file pinned-original-source.tex carries only \textbackslash{}input\{preamble\} on its first line, so no preamble text lives inside the chapter. The retained file sources/stacks-more-algebra/retained-preamble.tex is a mechanical, line-for-line copy of the project's single shared preamble as staged in the snapshot, namely sources/stacks-chow/preamble.tex lines 14-18 (Xy-pic setup) and 154-217 (theorem declarations and mathematical macros). Nothing was added, removed or redefined, no package was added, and the excerpt serves only as preamble, never as a span coordinate.\par\medskip
\noindent\textit{Editorial note.} Both counted spans start at the \textbackslash{}begin line of their environment, so the statement and the author's proof are delivered whole; surrounding prose of the section is deliberately outside the counted spans and is carried as standard prerequisites instead.\par\medskip
\noindent\textit{Editorial note.} The citation fields (source\_id, source\_version, authors, source\_url, license.name, source\_title) are identical to those of the cards already published from this chapter, so that the batch selector treats them as one source instead of deferring them as a same-source clash.\par\medskip
\noindent\textit{Allowed standard prerequisites (editorial).}\begin{itemize}
\item Derivation of a ring, the Leibniz rule, and the formula for extending a derivation to a localisation. definition discussion of the same section, source lines 11700-11752, before the counted statement. Supplies the object D, the rule that controls D on products and the explicit formula used in part (2).
\item The I-adic completion of a ring as the limit of the quotients R over I\textasciicircum{}n. used at proof lines 11769-11773 for part (1). Allows the compatible maps D\_n on quotients to be assembled into a derivation of the completion.
\item A finite type extension is generated by finitely many elements, and elements of R' are sums of monomials in those generators with coefficients in R. used at proof lines 11778-11794 for part (3). Reduces the statement about all of R' to the finitely many generators, which is the step that makes the power of g sufficient.
\end{itemize}\medskip
\ifdefined\hypertarget\hypertarget{sourcepointer-1}{}\fi
\noindent\textit{Source reference.} Original lemma 07PE: focused-seeds/AGP/sources/stacks-snapshot-20261011/tags.txt, source lines 9794--9794. The lemma's own canonical tag, quoted so that readers can resolve the label inside the delivered text; it points at the counted statement itself and carries no proof content.\par\medskip
\begin{lemma}
\label{lemma-derivation-extends}
Let $R$ be a ring. Let $D : R \to R$ be a derivation.
\begin{enumerate}
\item For any ideal $I \subset R$ the derivation $D$ extends
canonically to a derivation $D^\wedge : R^\wedge \to R^\wedge$
on the $I$-adic completion.
\item For any multiplicative subset $S \subset R$ the derivation
$D$ extends uniquely to the localization $S^{-1}R$ of $R$.
\end{enumerate}
If $R \subset R'$ is a finite type extension of rings such that
$R_g \cong R'_g$ for some $g \in R$ which is a nonzerodivisor in $R'$,
then $g^ND$ extends to $R'$ for some $N \geq 0$.
\end{lemma}

\begin{proof}
Proof of (1). For $n \geq 2$ we have $D(I^n) \subset I^{n - 1}$
by the Leibniz rule. Hence $D$ induces maps $D_n : R/I^n \to R/I^{n - 1}$.
Taking the limit we obtain $D^\wedge$. We omit the verification that
$D^\wedge$ is a derivation.

\medskip\noindent
Proof of (2). To extend $D$ to $S^{-1}R$ just set
$D(r/s) = D(r)/s - rD(s)/s^2$ and check the axioms.

\medskip\noindent
Proof of the final statement. Let $x_1, \ldots, x_n \in R'$ be generators
of $R'$ over $R$. Choose an $N$ such that $g^Nx_i \in R$.
Consider $g^{N + 1}D$. By (2) this extends to $R_g$. Moreover, by
the Leibniz rule and our construction of the extension above we have
$$
g^{N + 1}D(x_i) = g^{N + 1}D(g^{-N} g^Nx_i) = -Ng^Nx_iD(g) +
gD(g^Nx_i)
$$
and both terms are in $R$. This implies that
$$
g^{N + 1}D(x_1^{e_1} \ldots x_n^{e_n}) =
\sum e_i x_1^{e_1} \ldots x_i^{e_i - 1} \ldots x_n^{e_n} g^{N + 1}D(x_i)
$$
is an element of $R'$. Hence every element of $R'$ (which can be written
as a sum of monomials in the $x_i$ with coefficients in $R$) is mapped to an
element of $R'$ by $g^{N + 1}D$ and we win.
\end{proof}

\end{document}
